\documentclass[10pt,twocolumn]{article}

\usepackage[letterpaper,margin=0.72in]{geometry}
\usepackage[T1]{fontenc}
\usepackage[utf8]{inputenc}
\usepackage{lmodern}
\usepackage{microtype}
\usepackage{amsmath,amssymb}
\usepackage{array}
\usepackage{booktabs}
\usepackage{listings}
\usepackage{xcolor}
\usepackage{url}

\raggedbottom
\emergencystretch=2em

\lstdefinelanguage{EasyCrypt}{
  morekeywords={lemma,hoare,op,type,proof,qed,require,import,theory,end,by,conseq},
  sensitive=true,
  morecomment=[l]{//},
  morecomment=[s]{(*}{*)}
}

\lstset{
  language=EasyCrypt,
  basicstyle=\ttfamily\scriptsize,
  keywordstyle=\bfseries,
  columns=fullflexible,
  keepspaces=true,
  breaklines=true,
  frame=single,
  framerule=0.2pt,
  xleftmargin=2pt,
  xrightmargin=2pt,
  aboveskip=4pt,
  belowskip=4pt
}

\newcommand{\haetae}{HAETAE}
\newcommand{\ec}{EasyCrypt}
\newcommand{\jasmin}{Jasmin}
\newcommand{\Zq}{\mathbb{Z}_q}
\newcolumntype{P}[1]{>{\raggedright\arraybackslash}p{#1}}

\title{\vspace{-0.8em}Checked Functional Correctness of an Extracted \jasmin{} Scalar NTT for \haetae{}}
\author{Anonymous Submission}
\date{}

\begin{document}
\maketitle
\vspace{-1.1em}

\begin{abstract}
Number-theoretic transform (NTT) kernels have compact control structure, but
their correctness depends on low-level conventions that are easy to state
incorrectly: twiddle offsets, inverse scaling, Montgomery signs, and
intermediate bounds. We present a replayable \ec{} verification of the model
extracted from the scalar \haetae{} NTT routines written in \jasmin{}. The proof
connects exported procedure models to full 256-point NTT specifications through
field, Montgomery, algebraic, and refinement layers, while preserving the
implementation's table padding, inverse final scale, signed Montgomery
convention, and coefficient growth. The checked top-level results are two Hoare
triples over the extracted forward and inverse wrappers: from a 16-bit signed
coefficient representation, the forward wrapper computes the full NTT with a
24-bit output bound, while the inverse wrapper computes the specified inverse
transform in Montgomery form with a 16-bit output bound. The result is
deliberately kernel-level: it gives checked functional-correctness contracts for
the extracted scalar wrappers and makes the boundary to vectorized code, other
polynomial routines, full-scheme security, side-channel analysis, round-trip
composition, and fresh extraction regeneration explicit.
\end{abstract}

\section{Introduction}

Post-quantum signature schemes rely on polynomial arithmetic over finite rings.
In \haetae{}~\cite{HAETAE2022}, as in related lattice schemes, this arithmetic is
accelerated by a 256-point NTT modulo \(q=64513\). The implementation challenge
is that the routine is not merely a textbook transform. It is an in-place,
signed, Montgomery-domain computation with a bit-reversed twiddle schedule,
final inverse scaling, and machine-word representations. Small mistakes in
signs, table offsets, or intermediate bounds can invalidate a proof even when
tests pass. A kernel-level proof is therefore useful even without a full
signature-scheme proof: a higher-level proof can appeal to this scalar path only
after the transform, representation, and bit bounds have been tied to the
implementation model under explicit preconditions.

The verification described here targets the scalar \jasmin{} implementation,
with the main routines in \path{poly.jinc} and exported wrappers in
\path{hpoly.jazz}.
\jasmin{} is designed for high-assurance cryptographic implementations with
explicit control over low-level operations~\cite{Jasmin2017}. The proof is
carried out in \ec{}~\cite{EasyCrypt2013} over the extracted model and
culminates in two Hoare statements in \path{NTTEndToEnd.ec}. The
contribution is not a new NTT algorithm, but a checked correspondence between a
low-level extracted model and a mathematical specification that uses the same
constants and table schedule as the checked extracted routines and tables. The
distinction matters for this kernel: the forward implementation skips table
index 0, the inverse table reserves index 255 for final scaling, the Montgomery
constant is
\(q^{-1} \bmod 2^{32}\) rather than \(-q^{-1}\), and the checked forward bound is
24 bits rather than a smaller informal estimate.

This paper makes three contributions. First, it gives checked end-to-end
contracts for the extracted exported \jasmin{} forward and inverse NTT wrappers;
the exact theorem statements appear below and are Hoare triples over the
\path{Hpoly_extract.ec} procedure model. Second, it mechanizes the
\haetae{}-specific obligations that are usually hidden in informal NTT
arguments: the unused forward-table entry, the inverse table's separate final
scale, the \(q^{-1} \bmod 2^{32}\) Montgomery convention, and the \(16\to24\)
forward coefficient-bound growth. Third, it presents a proof organization that
keeps algebraic NTT correctness, executable loop correctness, word-array
representation, and extracted-code refinement separate until the final Hoare
composition.

The contribution is in this integration rather than in a new transform. The
checked statement reaches the \path{Hpoly_extract.ec} model of the public scalar
\jasmin{} entry points, not only an abstract NTT equation. The proof must align
the exact implementation schedule with the algebraic transform, because a
single off-by-one twiddle index or missing Montgomery factor changes the
postcondition. The artifact boundary is therefore part of the claim: the proof
establishes functional correctness for the extracted scalar wrappers, while
AVX2 code, other polynomial routines, side-channel behavior, full-scheme
correctness, and round-trip composition remain outside the displayed top-level
theorems.

This division also clarifies what is inherited from existing infrastructure and
what is specific to this work. \jasmin{} and \ec{} supply the programming and
proof environments. The artifact-specific work is the HAETAE instantiation:
the \(q=64513\) field and Montgomery constants, the concrete twiddle-table
agreement lemmas, the partial-transform algebra for the full 256-point NTT, and
the refinement invariants that preserve byte-array representation and bit
bounds through the extracted scalar program.

The acceptance criterion for the central claim is correspondingly concrete. A
reviewer should be able to trace the public \jasmin{} wrapper names to the
extracted program model, the extracted model to the imperative \ec{}
loops, and those loops to the full-transform specifications, while preserving
the stated representation and bit-bound predicates. The paper is organized
around that trace rather than around a standalone mathematical derivation of the
NTT.

\paragraph{Context and positioning.}
The work fits into computer-aided cryptography: \jasmin{} provides a low-level
language for high-assurance cryptographic code~\cite{Jasmin2017}, while
\ec{} supports mechanized reasoning about cryptographic programs and
specifications~\cite{EasyCrypt2013}. Broader work argues for tool-supported
links between mathematical claims and executable cryptographic
code~\cite{ComputerAidedCrypto2021}. This paper contributes a target-specific
instance of that methodology for the \haetae{} NTT. The semantic foundations are
provided by existing tools; the proof value is in closing the HAETAE-specific
gap between the extracted module and the EasyCrypt theorems, while
preserving constants, consumed table indices, representation predicates, and
coefficient bounds from the scalar implementation artifacts. For reviewers, the
intended claim is therefore narrow but
replayable: the extracted scalar wrappers satisfy the two displayed Hoare
contracts under their stated preconditions.

\section{Verification Problem}

The proof obligation is not simply that the textbook NTT equations are true.
The artifact must show that the concrete scalar program model implements the
specific full-transform maps used by the \ec{} specification, while preserving
the representation invariants required by later code. In this setting, a
seemingly small mismatch changes the theorem being proved. Table~\ref{tab:risks}
summarizes the implementation details that drive the verification task.

\begin{table}[t]
\centering
\scriptsize
\setlength{\tabcolsep}{3pt}
\begin{tabular}{P{0.26\columnwidth}P{0.36\columnwidth}P{0.28\columnwidth}}
\toprule
Detail & Why it matters & Representative evidence \\
\midrule
Forward table index 0 is padding & A full-array equality with mathematical
\texttt{zetas} is false; only indices \(1..255\) are consumed
& \path{NTT_Fq.ec}: \texttt{jzetas\_poly\_vals} \\
Inverse index 255 is scale & It is not a butterfly factor, so the invariant and
postcondition must isolate final scaling
& \path{NTT_Fq.ec}: \texttt{jzetas\_inv}\newline\texttt{\_255\_scaleE} \\
\texttt{QINV} sign convention & The reduction uses \(q^{-1}\), and the REDC
postcondition includes \(R^{-1}\)
& \path{Fq.ec}, \path{Montgomery.ec}: \texttt{SREDCp\_corr} \\
Unreduced butterfly sums & Coefficients grow during the forward transform, so
bit bounds are semantic postconditions
& \path{RefJasminNTT.ec}, \path{NTTEndToEnd.ec} \\
Byte-array representation & The public routines operate on 1024-byte arrays,
not directly on field arrays
& \path{NTT_Fq.ec}: \texttt{poly\_repr\_bound} \\
\bottomrule
\end{tabular}
\caption{Concrete details that the proof must account for.}
\label{tab:risks}
\end{table}

The inputs to the top-level Hoare triples are therefore not arbitrary
mathematical polynomials alone. They are byte-array values satisfying
\texttt{poly\_repr\_bound}: each 32-bit word decodes to the intended field
coefficient and its signed representative stays within the stated bit bound.
The outputs are specified in the same representation. This is what lets the
proof say both that the mathematical transform is computed and that the result
remains in a range compatible with the implementation-level arithmetic model.

The nontriviality is in making these obligations compose. The algebraic proof
is most naturally phrased over field arrays and roots of unity; the extracted
program operates on byte arrays, signed words, and Montgomery products; and the
actual loop schedule consumes only selected table intervals. The proof therefore
has to establish more than a high-level NTT identity: it must show that every
representation conversion, table lookup, Montgomery multiplication, and
coefficient-bound step preserves the exact preconditions needed by the next
layer.

\section{Implementation Target}

The forward routine modeled from \path{_poly_ntt} is an in-place Cooley--Tukey
style loop nest. It initializes \texttt{len} to 128 and halves it at each stage.
Inside each group it increments \texttt{zetasctr} before loading the twiddle, so
the forward table entry at index 0 is padding and the implementation consumes
\texttt{jzetas[1]} through \texttt{jzetas[255]}. Each butterfly computes a
Montgomery product \(t=\texttt{\_\_fqmul}(zeta, a[j+len])\), then writes
\(a[j]+t\) and \(a[j]-t\). No modular reduction is inserted between these
additions and subtractions; this is why coefficient bounds are part of the
functional postcondition.

The inverse routine modeled from \path{_poly_invntt} runs the inverse loop
direction, starting with \texttt{len = 1} and doubling it. Its butterfly table is
not indexed like the forward table. The implementation reads
\texttt{jzetas\_inv[0..254]} for the 255 inverse butterflies. After the loop it
performs a separate final pass using
\texttt{jzetas\_inv[255]}, whose value is the Montgomery-form scale factor
\(-29720\), satisfying the checked identity used to model multiplication by
\(256^{-1}\) together with the Montgomery convention.

The field-arithmetic component used for the extracted model comes from
\path{reduce.jinc}, which implements Montgomery
reduction~\cite{Montgomery1985}. The constant
\texttt{QINV = 940508161} is \(q^{-1} \bmod 2^{32}\), not \(-q^{-1}\). The
signed Montgomery reduction mirrors the implementation formula
\[
  (a - t q) / 2^{32}, \qquad
  t = \operatorname{int32}(a)\cdot \texttt{QINV},
\]
and the checked \ec{} theorem states the Montgomery result, namely multiplication
by \(R^{-1}\) modulo \(q\) for \(R=2^{32}\), together with the signed output
bound.

The object verified in \ec{} is the extracted procedure model
\path{Hpoly_extract.ec}. The top-level theorems do not reason directly over
\jasmin{} syntax; they reason over exported procedures and internal cores in
that extracted model. The refinement layer then connects those procedures to
the executable \ec{} NTT loops. This distinction is important for the claim:
source-level names identify the implementation target, while the checked Hoare
triples are statements about the extracted model and its imported semantics.

\begin{table}[t]
\centering
\scriptsize
\setlength{\tabcolsep}{3pt}
\begin{tabular}{P{0.29\columnwidth}P{0.31\columnwidth}P{0.30\columnwidth}}
\toprule
Layer & Main artifacts & Role in the proof \\
\midrule
\jasmin{} code & \path{hpoly.jazz}, \path{poly.jinc}, \path{reduce.jinc}, \path{zetas.jinc}
& Source of the extracted scalar implementation, tables, and Montgomery arithmetic \\
Checked extraction & \path{Hpoly_extract.ec}
& Extracted \ec{} model of exported wrappers and cores \\
Field and words & \path{Fq.ec}, \path{GFq.ec}, \path{BArray1024.ec}
& Signed residues, word arrays, field operations \\
Montgomery & \path{Montgomery.ec}
& REDC correctness with \(R^{-1}\) and signed bounds \\
Functional NTT & \path{NTT_Fq.ec}, \path{NTTFullSpec.ec}
& Executable loops and mathematical full-transform specs \\
Algebra & \path{NTTFullAlgebra.ec}
& Forward/inverse loop algebra and partial-transform invariants \\
Refinement & \path{RefJasminNTT.ec}
& Wrapper-to-core equivalence and core-to-loop refinement \\
End-to-end & \path{NTTEndToEnd.ec}
& Exported Hoare theorems for \path{poly_ntt_jazz} and \path{poly_invntt_jazz} \\
\bottomrule
\end{tabular}
\caption{Implementation artifacts and proof layers used by the checked result.}
\label{tab:layers}
\end{table}

\section{Specification and Algebra}

The \ec{} development separates executable loop models from mathematical full
transform specifications. The specifications operate on length-256 arrays over
\(\Zq\), with \(q=64513\), and use the checked root
\(\zeta=\texttt{zroot}=426\). In the specification, \(\zeta^{256}=-1\), and the
forward and inverse maps in \path{NTTFullSpec.ec} are defined by
\[
  \texttt{full\_ntt}(p)[i] =
    \sum_{j=0}^{255} p[j]\zeta^{(2\operatorname{br}(i)+1)j}
\]
and
\[
  \texttt{full\_invntt}(p)[i] =
    256^{-1}\sum_{j=0}^{255}
      p[j]\zeta^{-(2\operatorname{br}(j)+1)i},
\]
where \(\operatorname{br}\) is 8-bit reversal. The implementation-facing
procedures in \path{NTT_Fq.ec} then give executable loop models over field
coefficients, plus the representation predicate
\texttt{poly\_repr\_bound}. This predicate connects the byte-array
representation used by extracted \jasmin{} code to a polynomial over \(\Zq\),
while also recording a signed bit bound for each coefficient.

The bit reversal in the formulas is therefore part of the checked full-transform
specification, not only an artifact of loop order.

The forward twiddle array is defined mathematically as
\[
  \texttt{zetas}[k] = \zeta^{\operatorname{bsrev}_8(k)} \quad (0 \leq k < 256),
\]
where \(\operatorname{bsrev}_8\) is 8-bit reversal. At stage \(s\), group \(g\)
uses table index \(k=2^s+g\), matching the implementation's increment-before-use
counter. The checked table lemma for \texttt{jzetas} is intentionally restricted
to indices \(1\leq i<256\): index 0 is the implementation padding value 0, while
the mathematical Montgomery image of \texttt{zetas[0]} is \(R\). Thus the table
agreement claim is a Montgomery-form, consumed-index statement rather than a
raw equality between complete arrays.

The algebraic layer in \path{NTTFullAlgebra.ec} proves that the executable loops
implement the full specifications in \path{NTTFullSpec.ec}. A key technical
point is that the proof does not rely on a loose informal ``butterfly applies a
root'' argument. Instead it uses segment-indexed partial-transform definitions
and split lemmas for the low and high halves of each butterfly segment. This
matches the nested loop structure of the implementation: the invariant knows the
current length, segment start, and twiddle position, and states exactly which
portion of the transform has been established.

The inverse proof follows the same discipline but with the inverse loop order
and the final scale. The scale is not another butterfly twiddle: it is isolated
as index 255 of \texttt{jzetas\_inv}. The executable inverse procedure is
therefore specified as the inverse NTT followed by the Montgomery-form scaling
represented in \ec{} by \texttt{array256\_mont}.

\section{Mechanized Proof Strategy}

The mechanized development is organized so that each layer changes exactly one
view of the program. This is important for auditability: a reviewer can inspect
whether the algebra, the imperative loops, or the low-level representation is
responsible for a given claim. The procedure names in the following chain are the
checked \path{Hpoly_extract.ec} models of the exported \jasmin{} code and its
internal core. The forward path has the following shape:
\begin{center}
\scriptsize
\begin{tabular}{c}
\texttt{poly\_ntt\_jazz} \\
\(\Downarrow\) \texttt{poly\_ntt\_jazz\_ref} \\
\texttt{\_poly\_ntt} \\
\(\Downarrow\) \texttt{poly\_ntt\_core\_ref} \\
\texttt{NTT.ntt} \\
\(\Downarrow\) \texttt{ntt\_full\_ntt} \\
\texttt{full\_ntt}
\end{tabular}
\end{center}
The inverse path is analogous, using the corresponding inverse wrapper,
core-refinement, and algebra lemmas. The final file \path{NTTEndToEnd.ec}
composes these facts into the public Hoare triples by consequence.

Each arrow in this chain is a checked lemma with a narrower interface than the
whole theorem. The wrapper lemmas remove the public-call shell, the core
refinement lemmas connect extracted code to executable \ec{} loops, and the
algebra lemmas connect those loops to the full-transform specifications. This
decomposition prevents the end-to-end theorem from hiding either a
representation conversion or an algebraic NTT step inside an informal
equivalence.

The algebraic step is where the transform-specific proof work resides. The
specification file defines the forward and inverse maps as array-valued sums
over \(\Zq\). The algebra file proves that the imperative loops compute those
sums using stage invariants over partial transforms. For the forward transform,
the key split lemmas decompose a length-\(2^{k+1}\) partial sum into the two
length-\(2^k\) halves combined by a butterfly. For the inverse transform, the
proof follows the reverse loop order and accounts for the final multiplication
by the scale entry. This structure is why the proof can talk about the
implementation's nested loops without rewriting the top-level theorem in terms
of machine words.

The refinement step is deliberately separate. It relates the extracted
\jasmin{} procedures to the coefficient-array loop model while tracking
representation and bit bounds. The wrapper lemmas are simple, because the
public wrappers inline to the core routines. The core refinement lemmas are
larger: they maintain \texttt{poly\_repr\_bound}, prove that each loaded twiddle
is the intended field element on the actually consumed index range, and use the
signed Montgomery theorem to justify each call to the field multiplication
routine.

\section{End-to-End Theorems}

The final EasyCrypt file \path{NTTEndToEnd.ec} composes the algebraic loop
correctness theorems with refinement from the extracted procedure model. The two
exported results are the following Hoare triples, quoted in the same shape as the
checked file. Informally, the first theorem says that the extracted exported
forward procedure implements the mathematical full NTT and tracks the necessary
coefficient growth; the second says that the extracted exported inverse
procedure implements the mathematical inverse NTT in Montgomery representation
and restores a 16-bit bound.

\begin{lstlisting}
lemma haetae_jasmin_ntt_correct p :
  hoare [Hpoly_extract.M.poly_ntt_jazz :
    NTT_Fq.poly_repr_bound rp p 16 ==>
    NTT_Fq.poly_repr_bound res
      (NTTFullSpec.full_ntt p) 24].
\end{lstlisting}

\begin{lstlisting}
lemma haetae_jasmin_invntt_correct p :
  hoare [Hpoly_extract.M.poly_invntt_jazz :
    NTT_Fq.poly_repr_bound rp p 16 ==>
    NTT_Fq.poly_repr_bound res
      (NTT_Fq.array256_mont
        (NTTFullSpec.full_invntt p)) 16].
\end{lstlisting}

In these Hoare triples, \texttt{rp} is the byte-array input to the extracted
procedure model and \texttt{res} is the returned byte array. The predicate
\texttt{poly\_repr\_bound a p b} asserts both representation agreement with the
abstract polynomial \(p\) and a signed \(b\)-bit bound on the decoded
coefficients. Thus the postconditions combine a functional statement about the
computed transform with a range statement about the representation left for
later code. The inverse theorem deliberately postconditions
\texttt{array256\_mont (full\_invntt p)} rather than \texttt{full\_invntt p}:
the verified routine returns the Montgomery-form result specified by the
checked model.

The bit indices in these predicates have their ordinary signed-word meaning in
the development: \texttt{bw32 a sz} is the interval
\(-2^{sz} \leq \texttt{to\_sint}(a) < 2^{sz}\). Thus the forward theorem's
\(16\to24\) statement says that unreduced butterfly arithmetic is allowed to
grow but remains within the checked 24-bit signed range. The inverse theorem's
\(16\to16\) statement says that the inverse loop and final Montgomery-form scale
return to the checked 16-bit signed range for the Montgomery representation in
the postcondition.

The forward theorem says that, for any abstract polynomial \(p\), if the input
byte array \texttt{rp} represents \(p\) with signed coefficients bounded by 16
bits, then the extracted exported routine returns a byte array representing
\(\texttt{full\_ntt}(p)\) with a 24-bit bound. The inverse theorem has a
different but equally important shape: from a 16-bit represented input, it
returns the Montgomery representation of \(\texttt{full\_invntt}(p)\), again
bounded by 16 bits.

These bounds are part of the theorem, not post-hoc measurements. The forward
proof tracks growth through the unreduced additions and subtractions of the
butterfly network and ends at 24 bits. The inverse proof combines its loop
invariant with the final Montgomery scale to recover a 16-bit bound. This
distinction matters: a false forward claim such as \(16\to 23\) or an inverse
claim that ignores Montgomery scaling would not match the checked statements.

\section{Refinement Argument}

The refinement layer proves that the extracted \jasmin{} procedures behave like
the \ec{} loop models. In \path{RefJasminNTT.ec}, the wrapper lemmas
\path{poly_ntt_jazz_ref} and \path{poly_invntt_jazz_ref} show that the exported
wrappers call the internal extracted cores. The larger core lemmas
\path{poly_ntt_core_ref} and \path{poly_invntt_core_ref} relate those cores to
the executable \ec{} NTT procedures. These proofs must account for the
representation boundary:
\jasmin{} uses 32-bit words and 1024-byte arrays, while the functional
specification uses 256 field coefficients. Lemmas in \path{NTT_Fq.ec} bridge
these views through
\texttt{barray256\_to\_poly}, \texttt{word\_to\_coeff}, and
\texttt{poly\_repr\_bound}.

Montgomery multiplication is another boundary. The low-level code multiplies
signed 32-bit values in 64-bit registers and applies signed Montgomery reduction.
The \ec{} model clones a signed-reduction theory with \(q=64513\),
\(R=2^{32}\), \(q^{-1}=940508161\), and \(R^{-1}=50386\). Its core correctness
lemma proves that signed REDC returns a value congruent to \(aR^{-1}\) modulo
\(q\) and within the signed range used by the rest of the proof. This is the
reason the NTT proof can reason in \(\Zq\) while the implementation remains close
to the C and \jasmin{} arithmetic.

\section{Proof Obligations and Evidence}

The proof effort is concentrated around three obligations that correspond to
common failure modes in low-level NTT arguments. The first is representation
agreement. A coefficient of the mathematical polynomial is an element of
\(\Zq\), but the implementation stores a signed 32-bit representative in a
1024-byte array. The predicate \texttt{poly\_repr\_bound} bundles these two
facts: each word decodes to the corresponding field coefficient and its signed
integer representative is below a stated bit bound. The top-level theorems
therefore state both the transformed polynomial and the output size.

The second obligation is table agreement. The forward table cannot be stated as
a full-array equality between \texttt{jzetas} and the Montgomery image of
\texttt{zetas}, because index 0 differs by design. The checked lemmas therefore
prove the equality only on the consumed interval \(1\leq i<256\). The inverse
table has a different shape: indices \(0\) through \(254\) are butterfly factors,
and index \(255\) is reserved for the final scale. Treating this final entry as
another butterfly factor would give the wrong program invariant.

The third obligation is arithmetic agreement. The Montgomery-reduction model is
signed: it sign-extends 32-bit inputs, multiplies in 64-bit registers, and uses
an arithmetic right shift. The \ec{} proof connects that behavior to field
multiplication by proving the signed REDC theorem with the \(R^{-1}\) factor.
Consequently the field-level proof uses the same Montgomery convention as the
implementation, and the range conclusion supplies the fact needed by the NTT
bound invariants.

\begin{table}[t]
\centering
\scriptsize
\setlength{\tabcolsep}{3pt}
\begin{tabular}{P{0.31\columnwidth}P{0.25\columnwidth}P{0.34\columnwidth}}
\toprule
Claim checked & Boundary crossed & Representative evidence \\
\midrule
Exported forward wrapper satisfies final contract & Extracted wrapper to spec
& \path{NTTEndToEnd.ec}: \texttt{haetae\_jasmin}\newline\texttt{\_ntt\_correct} \\
Exported inverse wrapper satisfies final contract & Extracted wrapper to spec
& \path{NTTEndToEnd.ec}: \texttt{haetae\_jasmin}\newline\texttt{\_invntt\_correct} \\
Forward loop computes \texttt{full\_ntt} & Executable loop to math spec
& \path{NTTFullAlgebra.ec}: \texttt{ntt\_full\_ntt} \\
Inverse loop computes \texttt{full\_invntt} & Inverse loop plus final scale
& \path{NTTFullAlgebra.ec}: \texttt{invntt\_full\_invntt} \\
Forward wrapper equals core & Exported model to extracted core
& \path{RefJasminNTT.ec}: \texttt{poly\_ntt\_jazz\_ref} \\
Inverse wrapper equals core & Exported model to extracted core
& \path{RefJasminNTT.ec}: \texttt{poly\_invntt\_jazz\_ref} \\
Forward core refines loop model & Extracted core to \ec{} proc
& \path{RefJasminNTT.ec}: \texttt{poly\_ntt\_core\_ref} \\
Inverse core refines loop model & Extracted core to \ec{} proc
& \path{RefJasminNTT.ec}: \texttt{poly\_invntt\_core\_ref} \\
Consumed twiddle entries agree with specs & Concrete tables to field arrays
& \path{NTT\_Fq.ec}: \texttt{jzetas\_poly\_vals}, \texttt{jzetas\_inv\_255\_scaleE} \\
Signed Montgomery REDC is correct & Machine arithmetic to field arithmetic
& \path{Montgomery.ec}: \texttt{SREDCp\_corr} \\
\bottomrule
\end{tabular}
\caption{Examples of checked obligations used by the end-to-end result.}
\label{tab:evidence}
\end{table}

The reviewer-facing evidence is the theorem dependency chain ending in
\path{NTTEndToEnd.ec}: the end-to-end file imports the algebra and refinement
layers, composes their lemmas by consequence, and exposes the two Hoare triples
quoted above. The final file makes this composition explicit through small
intermediate lemmas: for each direction, one lemma combines the core-refinement
theorem with the algebraic loop theorem, and a second lemma moves from the
internal core to the exported wrapper. The displayed final theorems are
consequences of these checked compositions, not independent restatements of the
specification.

The top-level claim is replayed by invoking \texttt{easycrypt compile} on
\path{NTTEndToEnd.ec} with the development and extraction libraries on the load
path. Successful compilation checks the imported dependency chain under the
configured \ec{} environment, including the extracted module used by the
refinement proof. This replay is evidence for the displayed Hoare triples, not
a performance evaluation and not an extraction-regeneration check. It
establishes that the field, Montgomery, table-agreement, algebraic loop, and
refinement layers compose to the two displayed contracts for all inputs
satisfying the 16-bit \texttt{poly\_repr\_bound} precondition. A mismatch between
the prose and the named Hoare triples, table lemmas, or Montgomery lemma would
be an error in the paper, not an alternative interpretation of the verified
result.

\section{Discussion}

The development separates the mathematical transform, the executable loop, and
the low-level program until the final composition step. This organization keeps
changes in one layer from rewriting the obligations of the others. The
mathematical specification describes \texttt{full\_ntt} and
\texttt{full\_invntt}; the executable \ec{} loops expose the implementation's
control structure; and the refinement layer handles word-level details without
reopening the algebra.

This separation also makes several excluded claims visible. The checked
development does not state that the entire forward table equals the Montgomery
image of the mathematical \texttt{zetas} array, because that statement is false
at index 0. It does not state that Montgomery reduction preserves the input
residue without an \(R^{-1}\) factor, because that would be the wrong convention.
It does not state that the forward transform maps 16-bit inputs to 23-bit
outputs, because the checked postcondition is 24 bits. Each excluded statement
changes the shape of a proof obligation and would break the connection to the
checked artifact.

Compared with test-based validation of the NTT, the proof covers the full
preconditioned input space of the displayed Hoare triples. This matters for
cryptographic arithmetic because boundary cases arise near signed
representative limits and table offsets. The proof is still not a replacement
for all implementation validation: it assumes the semantics provided by the
\jasmin{} extraction and \ec{} libraries, and it targets the scalar extraction
used by the checked development. Constant-time behavior, compiler back-end
issues, and integration with the full signature scheme are separate questions.

The proof architecture suggests a practical pattern for other lattice kernels
with similar structure. First, fix the arithmetic convention, including the
Montgomery factor, before proving loop properties. Second, prove table agreement
only on the indices consumed by the implementation. Third, express loop progress
with partial-transform invariants that follow the concrete loop nest. Finally,
compose extracted-code refinement only after the algebraic loop theorem is
available. The \haetae{} proof is one instance of this pattern: it shows that
the low-level conventions are not side conditions to be checked informally, but
named obligations needed by the final Hoare contracts.

The scope is deliberately kernel-level. The result complements, rather than
replaces, full protocol or security proofs: it discharges a concrete
implementation-level arithmetic obligation for a performance-critical primitive.
This is the level at which table offsets, Montgomery conventions, and bit-growth
arguments are most likely to diverge from a high-level mathematical description.

\section{Related Work}

NTTs are finite-field analogues of FFT-style evaluation and interpolation
methods~\cite{Pollard1971}. They are now a standard implementation technique
for lattice-based cryptography, including schemes such as
CRYSTALS-Kyber~\cite{Kyber2018}, CRYSTALS-Dilithium~\cite{Dilithium2018}, and
\haetae{}~\cite{HAETAE2022}. This paper does not contribute a faster transform
or a new parameter set. Its subject is the narrower question of how to connect
one concrete scalar NTT implementation to a checked functional
specification without losing the implementation details that determine
correctness.

Verified cryptographic implementation projects have shown several ways to
connect executable arithmetic to formal specifications. HACL*~\cite{HACL2017}
verifies a broad C cryptographic library in the F*/Low* ecosystem, while
Fiat-Crypto~\cite{FiatCrypto2019} generates efficient finite-field arithmetic
from Coq-checked specifications. The present work is smaller in scope but
different in shape: it starts from hand-written low-level \jasmin{} routines,
uses the extracted \ec{} model, and proves that the resulting scalar NTT
wrappers satisfy HAETAE-specific transform contracts. It is therefore neither a
general verified library nor a code generator; it is a focused proof that a
concrete lattice-arithmetic kernel matches its specification despite low-level
representation, table-indexing, and Montgomery-domain details.

\ec{}~\cite{EasyCrypt2013} and the broader computer-aided cryptography
literature~\cite{ComputerAidedCrypto2021} provide the setting for mechanized
reasoning about cryptographic programs and claims. The present development is
not a game-based security proof for the \haetae{} signature scheme. It sits one
layer below such arguments: a functional-correctness proof for the concrete
arithmetic implementation that security proofs and algorithm descriptions would
otherwise have to trust informally. The proof therefore complements, rather
than competes with, protocol-level verification.

For NTTs and polynomial arithmetic, the closest comparison is the
\jasmin{}/\ec{} verification of Kyber/ML-KEM
implementations~\cite{KyberVerify2023}. That work is broader than the result
here because it connects Jasmin implementations for a KEM to \ec{}
specifications. The present paper is narrower and should be read as a kernel
proof, not as a full-scheme verification. Its technical delta is the
HAETAE-specific instantiation: modulus \(64513\), Montgomery constant
\(q^{-1} \bmod 2^{32}\), a full 256-point transform, a forward table with unused
padding at index 0, and an inverse table whose index 255 is final scaling rather
than a butterfly twiddle. These details are small enough to look like
engineering choices, but they determine the shape of the checked lemmas and
final Hoare postconditions. This is the main difference from a paper that only
states a textbook NTT proof or reports tests against reference vectors.

\section{Boundary and Trust}

The checked development establishes two functional-correctness contracts: one
for the scalar \jasmin{} forward NTT wrapper and one for the scalar inverse
wrapper. Both contracts are stated over the \ec{} model of the exported
procedures and the representation predicate used throughout the proof. This is
the intended theorem boundary: the result is a statement about extracted scalar
procedure models satisfying the two displayed Hoare triples.

Several adjacent claims are intentionally outside that boundary. The result does
not cover an AVX2 implementation, the full \haetae{} signature scheme,
side-channel security, or other exported polynomial routines such as
\path{poly_basemul_jazz}. It also does not expose a standalone checked
round-trip theorem stating that the verified forward routine followed by the
verified inverse routine is the identity on all represented inputs. Such a
statement would need to combine the two displayed Hoare triples with a checked
algebraic identity relating \texttt{full\_invntt} and \texttt{full\_ntt}, plus
the Montgomery-output convention of the inverse theorem. The current paper does
not depend on that missing composition theorem: the forward and inverse
contracts are stated and replayed separately.

The trusted base is correspondingly explicit. The proof relies on the semantics
of \ec{}, the \jasmin{} extraction used to obtain the \ec{} model of the source
procedures, the standard fixed-size word and array libraries imported by the
development, and the adequacy of the functional specifications
\texttt{full\_ntt} and \texttt{full\_invntt} as the intended \haetae{} NTT
contracts. Regenerating \path{Hpoly_extract.ec} from the \jasmin{} sources and
comparing it with the extraction used by the proof is an artifact-integrity
check, not part of the theorem replay reported in this paper. Validation of a
particular assembler, compiler backend, linker, or calling-context integration
is a separate engineering task.

\section{Conclusion}

This verification connects an extracted scalar polynomial-arithmetic model to a
precise formal specification without abstracting away the details that usually
cause NTT proof errors: signed Montgomery reduction, table padding,
inverse-table indexing, final scaling, and coefficient growth. The replayed
top-level EasyCrypt file composes table agreement, Montgomery correctness,
algebraic loop correctness, and refinement into two concrete Hoare contracts for
the extracted models of the exported \haetae{} NTT wrappers. The main lesson is methodological
as well as artifact-specific: for lattice-arithmetic kernels, the proof must
follow the implementation's representation, table schedule, and bounds closely
enough that small low-level conventions remain visible in the final theorem.

\bibliographystyle{abbrv}
\bibliography{refs}

\end{document}
